\documentclass[11pt,a4paper]{article}

\usepackage[margin=1in]{geometry}
\usepackage{amsmath,amssymb,mathtools,bm}
\usepackage{physics}
\usepackage{siunitx}
\usepackage{booktabs}
\usepackage{enumitem}
\usepackage{hyperref}
\usepackage{float}

\title{Complete Mathematical Formulation for Task-Space Tracking\\with Joint-Space Feasibility and CBF Safety}
\author{}
\date{}

\begin{document}
\maketitle
\tableofcontents
\newpage

%%%%%%%%%%%%%%%%%%%%%%%%%%%%%%%%%%%%%%%%%%%%%%%%%%%%%%%%%%%%
\section{Problem Statement and Why It Is Nontrivial}
%%%%%%%%%%%%%%%%%%%%%%%%%%%%%%%%%%%%%%%%%%%%%%%%%%%%%%%%%%%%

We are given an end-effector trajectory from a CSV file and want a 7-DOF robot arm to execute it.
The trajectory is specified in task space:
\[
\mathbf{p}_r(t)\in\mathbb{R}^3,\qquad R_r(t)\in SO(3)\ \text{(if orientation is available)}.
\]
However, the robot is actuated in joint space:
\[
\mathbf{q}(t)\in\mathbb{R}^7.
\]
Hence, each control step must transform a Cartesian pose-tracking objective into feasible joint velocities while satisfying hard physical limits and safety constraints.

The key difficulty is that these goals can conflict:
\begin{enumerate}[label=\arabic*.]
\item Follow the task-space path accurately.
\item Stay inside a safety region in task space (CBF constraint).
\item Respect joint position, velocity, and acceleration limits.
\item Use configuration redundancy (7 DOF) intelligently via secondary objectives.
\end{enumerate}

The final controller must produce commands that are both \emph{geometrically correct} (task-space tracking) and \emph{physically executable} (joint-space feasibility).

%%%%%%%%%%%%%%%%%%%%%%%%%%%%%%%%%%%%%%%%%%%%%%%%%%%%%%%%%%%%
\section{Kinematics and Variables}
%%%%%%%%%%%%%%%%%%%%%%%%%%%%%%%%%%%%%%%%%%%%%%%%%%%%%%%%%%%%

\subsection{Joint-space variables}
\[
\mathbf{q}\in\mathbb{R}^7,\qquad
\dot{\mathbf{q}}\in\mathbb{R}^7,\qquad
\ddot{\mathbf{q}}\in\mathbb{R}^7.
\]

\subsection{End-effector pose}
\[
\mathbf{p}=f_p(\mathbf{q})\in\mathbb{R}^3,\qquad
R=f_R(\mathbf{q})\in SO(3).
\]

\subsection{Differential kinematics}
\[
\dot{\mathbf{p}} = J_p(\mathbf{q})\,\dot{\mathbf{q}},\qquad
\bm{\omega}=J_\omega(\mathbf{q})\,\dot{\mathbf{q}},
\]
with
\[
J_p\in\mathbb{R}^{3\times 7},\quad
J_\omega\in\mathbb{R}^{3\times 7}.
\]

If orientation is tracked, define:
\[
J_6=\begin{bmatrix}J_p\\J_\omega\end{bmatrix}\in\mathbb{R}^{6\times 7},\qquad
\mathbf{v}_6=\begin{bmatrix}\dot{\mathbf{p}}\\\bm{\omega}\end{bmatrix}.
\]

%%%%%%%%%%%%%%%%%%%%%%%%%%%%%%%%%%%%%%%%%%%%%%%%%%%%%%%%%%%%
\section{Reference Construction from CSV}
%%%%%%%%%%%%%%%%%%%%%%%%%%%%%%%%%%%%%%%%%%%%%%%%%%%%%%%%%%%%

At sample time \(t_k\), position reference is \(\mathbf{p}_r(t_k)\). If orientation columns exist, we also have \(R_r(t_k)\).

\subsection{Feedforward translational velocity}
\[
\dot{\mathbf{p}}_r(t_k)\approx
\frac{\mathbf{p}_r(t_{k+1})-\mathbf{p}_r(t_k)}{t_{k+1}-t_k}.
\]

\subsection{Feedforward angular velocity}
\[
\bm{\omega}_r(t_k)\approx
\frac{\log\!\left(R_r(t_k)^\top R_r(t_{k+1})\right)^\vee}{t_{k+1}-t_k}.
\]

%%%%%%%%%%%%%%%%%%%%%%%%%%%%%%%%%%%%%%%%%%%%%%%%%%%%%%%%%%%%
\section{Tracking Errors and Desired Cartesian Motion}
%%%%%%%%%%%%%%%%%%%%%%%%%%%%%%%%%%%%%%%%%%%%%%%%%%%%%%%%%%%%

\subsection{Position error}
\[
\mathbf{e}_p=\mathbf{p}_r-\mathbf{p}.
\]

\subsection{Orientation error}
Define:
\[
R_e=R_rR^\top.
\]
Use the standard vector form:
\[
\mathbf{e}_\omega
=
\frac12
\begin{bmatrix}
R_{e32}-R_{e23}\\
R_{e13}-R_{e31}\\
R_{e21}-R_{e12}
\end{bmatrix}.
\]

\subsection{Desired Cartesian command}
\[
\mathbf{v}_{d,p}=\dot{\mathbf{p}}_r+K_p\mathbf{e}_p.
\]
If orientation is available:
\[
\mathbf{v}_{d,\omega}=\bm{\omega}_r+K_o\mathbf{e}_\omega,
\qquad
\mathbf{v}_{6,d}=
\begin{bmatrix}
\mathbf{v}_{d,p}\\
\mathbf{v}_{d,\omega}
\end{bmatrix}.
\]
If orientation is not available, only \((J_p,\mathbf{v}_{d,p})\) is used.

%%%%%%%%%%%%%%%%%%%%%%%%%%%%%%%%%%%%%%%%%%%%%%%%%%%%%%%%%%%%
\section{Task-Space Safety: Position CBF}
%%%%%%%%%%%%%%%%%%%%%%%%%%%%%%%%%%%%%%%%%%%%%%%%%%%%%%%%%%%%

Define a spherical safe set with center \(\mathbf{c}\) and radius \(X_m\):
\[
h_p(\mathbf{p})=X_m^2-\norm{\mathbf{p}-\mathbf{c}}^2.
\]
Safe region:
\[
\mathcal{C}=\{\mathbf{p}: h_p(\mathbf{p})\ge 0\}.
\]

Apply CBF condition:
\[
\dot{h}_p+\gamma_p h_p\ge 0,\qquad \gamma_p>0.
\]
Since
\[
\dot{h}_p=-2(\mathbf{p}-\mathbf{c})^\top \dot{\mathbf{p}}
       =-2(\mathbf{p}-\mathbf{c})^\top J_p\dot{\mathbf{q}},
\]
we get a linear inequality in \(\dot{\mathbf{q}}\):
\[
2(\mathbf{p}-\mathbf{c})^\top J_p\dot{\mathbf{q}}
\le
\gamma_p\left(X_m^2-\norm{\mathbf{p}-\mathbf{c}}^2\right).
\]

This is enforced as a hard QP constraint at each step.

\subsection{Velocity CBF Condition}
To enforce a dynamic speed safety layer, define
\[
h_v(\dot{\mathbf{p}})=V_m^2-\norm{\dot{\mathbf{p}}}^2.
\]
The safe set is
\[
\mathcal{C}_v=\{\dot{\mathbf{p}}: h_v\ge 0\}.
\]
Apply the CBF condition
\[
\dot{h}_v+\gamma_v h_v\ge 0,\qquad \gamma_v>0.
\]
Since
\[
\dot{h}_v=-2\dot{\mathbf{p}}^\top\ddot{\mathbf{p}},
\]
we get
\[
-2\dot{\mathbf{p}}^\top\ddot{\mathbf{p}}+\gamma_v\left(V_m^2-\norm{\dot{\mathbf{p}}}^2\right)\ge 0.
\]

For implementation with decision variable \(\dot{\mathbf{q}}_k\), use
\[
\dot{\mathbf{p}}_k \approx J_p\dot{\mathbf{q}}_{k-1},\qquad
\ddot{\mathbf{p}}_k \approx \frac{J_p\dot{\mathbf{q}}_k-\dot{\mathbf{p}}_k}{\Delta t_k}.
\]
Substitute into the CBF inequality:
\[
\dot{\mathbf{p}}_k^\top J_p\dot{\mathbf{q}}_k
\le
\norm{\dot{\mathbf{p}}_k}^2
\frac{\gamma_v\Delta t_k}{2}\left(V_m^2-\norm{\dot{\mathbf{p}}_k}^2\right).
\]
This is linear in \(\dot{\mathbf{q}}_k\) and can be directly added to \(A_{ineq}\dot{\mathbf{q}}\le \mathbf{b}_{ineq}\).

%%%%%%%%%%%%%%%%%%%%%%%%%%%%%%%%%%%%%%%%%%%%%%%%%%%%%%%%%%%%
\section{Secondary Joint-Space Objective}
%%%%%%%%%%%%%%%%%%%%%%%%%%%%%%%%%%%%%%%%%%%%%%%%%%%%%%%%%%%%

To exploit redundancy, define a preferred joint velocity:
\[
\dot{\mathbf{q}}_{sec}
=
k_h(\mathbf{q}_{home}-\mathbf{q})
-k_\ell \nabla V_{lim}(\mathbf{q}).
\]

Typical limit-avoidance potential:
\[
V_{lim}(\mathbf{q})=
\sum_{i=1}^7
\left(
-\log(q_i-q_i^{\min})-\log(q_i^{\max}-q_i)
\right),
\]
whose gradient grows near limits and pushes joints inward.

%%%%%%%%%%%%%%%%%%%%%%%%%%%%%%%%%%%%%%%%%%%%%%%%%%%%%%%%%%%%
\section{Updated Optimization Problem (Per Time Step)}
%%%%%%%%%%%%%%%%%%%%%%%%%%%%%%%%%%%%%%%%%%%%%%%%%%%%%%%%%%%%

At each \(k\), solve for \(\dot{\mathbf{q}}_k\):
\[
\min_{\dot{\mathbf{q}}}
\underbrace{\norm{J\dot{\mathbf{q}}-\mathbf{v}_d}_2^2}_{\text{task tracking}}
+\underbrace{\lambda^2\norm{\dot{\mathbf{q}}}_2^2}_{\text{damping}}
+\underbrace{w_{smooth}\norm{\dot{\mathbf{q}}-\dot{\mathbf{q}}_{k-1}}_2^2}_{\text{temporal smoothness}}
+\underbrace{w_{null}\norm{\dot{\mathbf{q}}-\dot{\mathbf{q}}_{sec}}_2^2}_{\text{secondary behavior}}.
\]

Here:
\[
J=
\begin{cases}
J_p & \text{position-only tracking}\\
J_6 & \text{position + orientation tracking}
\end{cases},
\qquad
\mathbf{v}_d=
\begin{cases}
\mathbf{v}_{d,p}\\
\mathbf{v}_{6,d}
\end{cases}.
\]

\subsection{Equivalent QP matrix form}
\[
\min_{\dot{\mathbf{q}}}\;
\frac12\dot{\mathbf{q}}^\top H\dot{\mathbf{q}}-g^\top\dot{\mathbf{q}},
\]
with
\[
H=J^\top J + \left(\lambda^2+w_{smooth}+w_{null}\right)I,
\]
\[
g=J^\top\mathbf{v}_d + w_{smooth}\dot{\mathbf{q}}_{k-1} + w_{null}\dot{\mathbf{q}}_{sec}.
\]

%%%%%%%%%%%%%%%%%%%%%%%%%%%%%%%%%%%%%%%%%%%%%%%%%%%%%%%%%%%%
\section{Hard Constraints in the QP}
%%%%%%%%%%%%%%%%%%%%%%%%%%%%%%%%%%%%%%%%%%%%%%%%%%%%%%%%%%%%

\subsection{Joint position feasibility (next-step form)}
With margin \(\delta\):
\[
\mathbf{q}_{min}+\delta \le \mathbf{q}_k+\dot{\mathbf{q}}\Delta t_k \le \mathbf{q}_{max}-\delta,
\]
which gives:
\[
\frac{\mathbf{q}_{min}+\delta-\mathbf{q}_k}{\Delta t_k}
\le \dot{\mathbf{q}} \le
\frac{\mathbf{q}_{max}-\delta-\mathbf{q}_k}{\Delta t_k}.
\]

\subsection{Joint velocity limits}
\[
-\dot{\mathbf{q}}_{max}\le \dot{\mathbf{q}}\le \dot{\mathbf{q}}_{max}.
\]

\subsection{Joint acceleration limits}
\[
\dot{\mathbf{q}}_{k-1}-\ddot{\mathbf{q}}_{max}\Delta t_k
\le
\dot{\mathbf{q}}
\le
\dot{\mathbf{q}}_{k-1}+\ddot{\mathbf{q}}_{max}\Delta t_k.
\]

All three are merged into box bounds:
\[
\mathbf{l}_b \le \dot{\mathbf{q}} \le \mathbf{u}_b.
\]

\subsection{Task-space CBF inequality}
\[
2(\mathbf{p}-\mathbf{c})^\top J_p\dot{\mathbf{q}}
\le
\gamma_p\left(X_m^2-\norm{\mathbf{p}-\mathbf{c}}^2\right).
\]

\subsection{Velocity CBF inequality}
Let \(\dot{\mathbf{p}}_k=J_p\dot{\mathbf{q}}_{k-1}\). Then impose
\[
\dot{\mathbf{p}}_k^\top J_p\dot{\mathbf{q}}_k
\le
\norm{\dot{\mathbf{p}}_k}^2
\frac{\gamma_v\Delta t_k}{2}\left(V_m^2-\norm{\dot{\mathbf{p}}_k}^2\right).
\]
This provides a barrier-style speed safety condition in addition to the norm bound.

\subsection{Task-space speed bound (conservative linearization)}
The Euclidean bound \(\norm{\dot{\mathbf{p}}}_2\le V_m\) is conservatively enforced via:
\[
\left|[J_p\dot{\mathbf{q}}]_j\right| \le \frac{V_m}{\sqrt{3}},\quad j=1,2,3.
\]

\subsection{Task-space acceleration bound (conservative discrete form)}
Using \(\dot{\mathbf{p}}\approx J_p\dot{\mathbf{q}}\):
\[
\left|
\frac{[J_p\dot{\mathbf{q}}]_j - [J_p\dot{\mathbf{q}}_{k-1}]_j}{\Delta t_k}
\right|
\le
\frac{A_m}{\sqrt{3}},\quad j=1,2,3.
\]

Thus all task constraints are linear inequalities:
\[
A_{ineq}\dot{\mathbf{q}}\le \mathbf{b}_{ineq}.
\]

%%%%%%%%%%%%%%%%%%%%%%%%%%%%%%%%%%%%%%%%%%%%%%%%%%%%%%%%%%%%
\section{Final QP Solved in Implementation}
%%%%%%%%%%%%%%%%%%%%%%%%%%%%%%%%%%%%%%%%%%%%%%%%%%%%%%%%%%%%

At each step:
\[
\begin{aligned}
\min_{\dot{\mathbf{q}}}\quad&
\frac12\dot{\mathbf{q}}^\top H\dot{\mathbf{q}}-g^\top\dot{\mathbf{q}}\\
\text{s.t.}\quad&
\mathbf{l}_b\le \dot{\mathbf{q}}\le \mathbf{u}_b,\\
&A_{ineq}\dot{\mathbf{q}}\le \mathbf{b}_{ineq}.
\end{aligned}
\]

This is solved using OSQP when available; otherwise a projected-QP fallback is used.

%%%%%%%%%%%%%%%%%%%%%%%%%%%%%%%%%%%%%%%%%%%%%%%%%%%%%%%%%%%%
\section{Discrete-Time State Update}
%%%%%%%%%%%%%%%%%%%%%%%%%%%%%%%%%%%%%%%%%%%%%%%%%%%%%%%%%%%%

After obtaining \(\dot{\mathbf{q}}_k\):
\[
\mathbf{q}_{k+1}=\mathbf{q}_k+\dot{\mathbf{q}}_k\Delta t_k.
\]
Numerical safeguard:
\[
\mathbf{q}_{k+1}\leftarrow
\operatorname{clip}(\mathbf{q}_{k+1},\mathbf{q}_{min},\mathbf{q}_{max}).
\]

%%%%%%%%%%%%%%%%%%%%%%%%%%%%%%%%%%%%%%%%%%%%%%%%%%%%%%%%%%%%
\section{Validation and Logged Quantities}
%%%%%%%%%%%%%%%%%%%%%%%%%%%%%%%%%%%%%%%%%%%%%%%%%%%%%%%%%%%%

Typical diagnostics:
\[
e_p(k)=\norm{\mathbf{p}_r(k)-\mathbf{p}(k)},
\qquad
e_o(k)=\norm{\mathbf{e}_\omega(k)}\ \text{(if orientation tracked)}.
\]

CBF residual:
\[
r_{cbf}(k)=\dot{h}_p(k)+\gamma_p h_p(k),
\]
with \(\dot{h}_p\) computed from executed task velocity.

Task bounds:
\[
\norm{\dot{\mathbf{p}}(k)}\le V_m,\qquad
\norm{\ddot{\mathbf{p}}(k)}\le A_m.
\]

Joint bounds:
\[
\mathbf{q}_{min}\le \mathbf{q}(k)\le \mathbf{q}_{max},\quad
\norm{\dot{\mathbf{q}}(k)}_\infty\le \dot{\mathbf{q}}_{max},\quad
\norm{\ddot{\mathbf{q}}(k)}_\infty\le \ddot{\mathbf{q}}_{max}.
\]

%%%%%%%%%%%%%%%%%%%%%%%%%%%%%%%%%%%%%%%%%%%%%%%%%%%%%%%%%%%%
\section{Algorithm (Implementation-Level)}
%%%%%%%%%%%%%%%%%%%%%%%%%%%%%%%%%%%%%%%%%%%%%%%%%%%%%%%%%%%%

At each sample \(k\):
\begin{enumerate}[label=\arabic*.]
\item Compute current pose, Jacobians, and errors.
\item Build \(\mathbf{v}_d\) (position only or full 6D with orientation).
\item Build \(H,g\) from tracking + smoothness + null preference.
\item Build hard bounds \(\mathbf{l}_b,\mathbf{u}_b\) from joint position/velocity/acceleration constraints.
\item Build linear inequalities for CBF, task velocity, and task acceleration.
\item Solve constrained QP for \(\dot{\mathbf{q}}_k\).
\item Integrate to \(\mathbf{q}_{k+1}\), clip numerically, log diagnostics.
\end{enumerate}

%%%%%%%%%%%%%%%%%%%%%%%%%%%%%%%%%%%%%%%%%%%%%%%%%%%%%%%%%%%%
\section{Interpretation}
%%%%%%%%%%%%%%%%%%%%%%%%%%%%%%%%%%%%%%%%%%%%%%%%%%%%%%%%%%%%

This formulation provides a clear priority structure:
\begin{itemize}
\item The objective tracks the desired Cartesian motion.
\item Hard constraints guarantee safety and feasibility.
\item Secondary terms shape posture without violating primary constraints.
\end{itemize}

In short: it is a mathematically consistent way to stay close to the desired task-space trajectory while guaranteeing joint-space executability and task-space safety.

%%%%%%%%%%%%%%%%%%%%%%%%%%%%%%%%%%%%%%%%%%%%%%%%%%%%%%%%%%%%
\section{How to Compile}
%%%%%%%%%%%%%%%%%%%%%%%%%%%%%%%%%%%%%%%%%%%%%%%%%%%%%%%%%%%%

Example:
\begin{verbatim}
pdflatex TaskJoint_QP_Formulation.tex
\end{verbatim}
Run twice if needed for table of contents.

\end{document}
